\chapter{The Two Faces at the Boundary}
% OUTLINE Ch8 "The Two Faces at the Boundary". THE HONESTY-CRITICAL
% CHAPTER. Laws: CITATION (Einstein 1915; Dirac 1930; Madelung 1927; the
% founding text of the apparatus cited once, per the outline row), the
% Ch3->Ch8 SEAM (the covariant/contravariant words point at Ch3's CITED
% construction, never device coinage), the FOUNDING-PROGRAM BLOCK stated
% AS PROMISE with its zero-energy self-fence, NEVER as theorem; NS owed;
% ends at the door. Beastmaster reads the founding-program block from
% source on arrival (pre-arranged); Kodo per-claim gate.

The language is assembled, every piece cited and every price public, and
this chapter does what the whole construction was for: it reads. Two
descriptions dominated the physics of the twentieth century. Both are
written in exactly the language of the last seven chapters, and --- the
fact this book has been walking toward --- each of them leans on a
different one of Chapter 3's two species. One is a covariant-face
description; the other is contravariant-face. Neither face was invented
for its theory; the reader watched both fall out of the chain rule,
cited, before either theory entered the room.

Read the geometric description first. In November 1915 Einstein published
the field equations of gravitation, and in this book's language the
sentence is short. Take the ruler --- in its Lorentzian signature, the one
paragraph of Chapter 7 --- and let it select its carrying rule by the
fundamental theorem; feed that rule to Chapter 6's loop-machine; contract
the resulting curvature to its Ricci reading and scalar, the smaller
readings Chapter 6 named. That assembly --- curvature, read down to two
slots --- stands on the equation's left. On its right stands a two-slot
machine holding the content of matter and field --- the stress-energy,
the direct heir of the machine Cauchy found inside loaded bodies in 1822.
The equation says the two machines agree, slot for slot, at every event:
geometry's curvature reads matter's content. And the description's whole
grammar runs on the lowered face: the metric, the Ricci reading, the
stress-energy --- two covector slots each, machines awaiting arrows, the
covariant species of Chapter 3 doing all the load-bearing work.

Now the quantum description. In 1930 Dirac published the book that gave
the theory its grammar, and the grammar is Chapter 3's other face. A
state of a system is a ket --- an arrow, not in the spacetime manifold
but in a state-space of its own, linear from the outset; an instrument's
question is a bra --- a read-out, an element of the dual; and every
prediction the theory makes is the elementary act this book met in
Chapter 3, instrument-meets-state, the bracket that gives the formalism
its name. The description's working face is the raised one: the state
side, the arrow species, contravariant under its changes of basis. One
honesty note belongs here, and Chapter 3 already paid for it: what the
two descriptions share is the \emph{grammar} --- the two species and
their pairing --- not the stage. The quantum arena is not the spacetime
manifold, and this book does not blur the two; it observes only that
both theories speak the constructed language, each favoring one face.

Between the two faces stands one more cited result, and it is the
closest thing the mathematics owns to a stitch. In 1927 --- before
Dirac's book --- Erwin Madelung showed that the Schr\"odinger equation
for the electron can be rewritten, exactly, as a pair of fluid
equations: a continuity equation for a density, and a momentum equation
of Euler's form carrying one additional term the trade calls the quantum
potential. Not an analogy; a change of variables. The electron's quantum
description genuinely has a hydrodynamic form, and the reader of this
series' preface will recognize what that means: the fourth station ---
the fluid description, the sounding not yet taken --- has, in Madelung,
its nearest earned neighbor. What separates Madelung's fluid from the
fluid of the founding promise is, among other things, viscosity: the
dissipative structure of the Navier--Stokes equations is not in the
1927 result, and this book, which has paid every debt in the open,
marks this one still owed.

\begin{fence}
The founding program, stated as what it is. The founding text of the
apparatus this series documents --- the same opening pages that owed the
fluid its debt --- sets the series its deepest assignment in one
sentence: that the covariance of the geometric description will never
settle, at the boundary, with the contravariance of the quantum one, no
matter how much energy is there. And the text fences itself in the next
breath, allowing the empty case: at zero energy there is nothing to
disagree about. This book states that sentence as what it is --- a
\emph{promise} the series' reader is to be shown, made by the founding
text, in words --- covariance, contravariance --- that Chapter 3 of this
book constructed from the chain rule and cited to their owners. It is
not a theorem. No volume of this series has proved it, this one least
of all; the two readings above show only that both descriptions speak
one language while favoring opposite faces --- which is what makes the
boundary sayable, and nothing about what happens there. Whether the
promise can be earned, and in which volume's terms, stays on the far
side of the same fence that holds the fluid: owed, named, and not
crossed here.
\end{fence}

And that is the door at which this book ends. It was never this volume's
work to settle a boundary; its work was to build, from cited topology
and cited linear algebra, a language in which charts, arrows, machines,
carrying, curvature, and measure all stand with their prices published
--- and then to demonstrate the language by reading the century's two
great descriptions in it, one on each face. The reader now holds what
the preface promised: the ground the four soundings stand on. Three of
the four stations have appeared in these pages --- the geometric, the
quantum, and, at one remove, the electron model whose measured number
this series' first volume brackets; the fourth is owed, and the door it
waits behind is now, at least, a door with a frame. What remains for
this volume is only its accounting: the ledger of names and dates this
book has spent, chapter by chapter, which is the back matter's business
and the next and final chapter.
